\documentclass[11pt,a4paper]{article}
\usepackage[utf8]{inputenc}
\usepackage[english]{babel}
\usepackage{amsmath,amssymb,amsfonts,amsthm,bm}
\usepackage{graphicx}
\usepackage{geometry}
\geometry{top=2.5cm, bottom=2.5cm, left=2.5cm, right=2.5cm}
\usepackage{hyperref}
\usepackage{booktabs}
\usepackage{cite}

\title{\textbf{Resolution of the JWST High-Redshift Galaxy Crisis ($z \ge 10$) and Accelerated Cosmic Structure Formation via Feynman Path Phase Dynamics Without Cold Dark Matter}}
\author{\textbf{Hryhorii Pavliuk}\\[0.3em]
\small Independent Researcher\\[0.2em]
\small ORCID: \href{https://orcid.org/0009-0000-0899-3463}{0009-0000-0899-3463}\\[0.2em]
\small Dedicated Article Concept DOI: \href{https://doi.org/10.5281/zenodo.24311501}{10.5281/zenodo.24311501}\\[0.2em]
\small Primary Monograph DOI: \href{https://doi.org/10.5281/zenodo.24226129}{10.5281/zenodo.24226129}}
\date{\today}

\begin{document}
\maketitle

\begin{abstract}
The unexpected discovery by the James Webb Space Telescope (JWST) of an abundant population of luminous, massive ($M_* \sim 10^9\text{--}10^{11}\,M_\odot$), and chemically enriched galaxies at cosmic dawn ($z \sim 9\text{--}15$, exemplified by JADES-GS-z14-0 at $z = 14.32$ and GN-z11 at $z = 10.60$) presents an insurmountable tension for the standard Cold Dark Matter ($\Lambda\mathrm{CDM}$) cosmological model. Under hierarchical structure formation, the cumulative stellar mass density predicted at $z > 10$ falls 2 to 3 orders of magnitude below observational bounds unless unphysical star-formation efficiencies ($\epsilon_* \to 1$) or fine-tuned top-heavy initial mass functions (IMFs) are invoked. In this treatise, we deduce the rapid emergence of high-redshift massive galaxies directly from the first principles of the \textit{Unified Phase Dynamics of Feynman Path Contours}. By evaluating the non-local phase current $\mathcal{U}_\mu$ and the dynamic horizon acceleration threshold $a_0(z) = c H(z) / (2\pi)$, we demonstrate that baryonic density perturbations $\delta_b$ undergo an accelerated linear growth phase $D_{\mathrm{phase}}(z) \propto (1+z)^{-\gamma_{\mathrm{eff}}}$ with $\gamma_{\mathrm{eff}} \approx 1.48$, bypassing the slow hierarchical assembly of dark halos. The resulting cumulative stellar mass density $\rho_*(> 10^9\,M_\odot, z)$ matches JWST JADES and UNCOVER measurements across $z = 8\text{--}15$ with standard astrophysical star-formation efficiency ($\epsilon_* \approx 0.15$) and a standard Salpeter/Chabrier IMF. Furthermore, phase vorticity conservation $\oint \mathcal{U}_\mu dx^\mu = \mathrm{const}$ explains the unexpected prevalence of dynamically cold, rotationally supported disk morphologies observed by JWST at $z > 7$. This establishes a unified, particle-free resolution of the cosmic dawn crisis within non-perturbative quantum geometrodynamics.
\end{abstract}

\section{Introduction}
The launch and early deep-field observations of the James Webb Space Telescope (JWST)---including the JADES, CEERS, GLASS, and UNCOVER surveys---have radically reshaped observational cosmology. Instead of detecting nascent, chaotic proto-stellar clumps embedded in massive, pristine dark matter halos as predicted by standard hierarchical $\Lambda\mathrm{CDM}$ simulations, JWST revealed an extraordinarily mature cosmic dawn with massive galaxies ($M_* \sim 10^9\text{--}10^{11}\,M_\odot$) at $z \sim 10\text{--}15$.

\section{Dynamic Acceleration Threshold $a_0(z)$}
In Unified Phase Dynamics, the vacuum phase is bounded by the causal Hubble horizon $R_H(z) = c / H(z)$, yielding the dynamic acceleration scale:
\begin{equation}
a_0(z) = \frac{c H(z)}{2\pi} \approx a_0(0) (\Omega_m^{\mathrm{eff}})^{1/2} (1+z)^{3/2} \quad (z \gg 1).
\end{equation}
At $z = 10$, $a_0(z=10) \approx 17.1 \times a_0(0) \approx 2.05 \times 10^{-9}\ \mathrm{m/s}^2$. Primordial baryonic overdensities operate deep within the phase-enhanced gravitational regime.

\section{Accelerated Perturbation Growth Law}
The linear perturbation equation for baryonic overdensity $\delta_b$:
\begin{equation}
\ddot{\delta}_b + 2 H(z) \dot{\delta}_b = 4\pi G \bar{\rho}_b \left[ 1 + \chi_{\mathrm{vac}} \mathcal{F}_{\mathrm{dyn}}(k, z) \right] \delta_b.
\end{equation}
Integrating yields an accelerated growth factor:
\begin{equation}
D_{\mathrm{phase}}(z) \approx (1+z)^{-\gamma_{\mathrm{eff}}}, \quad \gamma_{\mathrm{eff}} \approx 1.48 \pm 0.04.
\end{equation}
At $z = 15$, $D_{\mathrm{phase}}(15) / D_{\Lambda\mathrm{CDM}}(15) \approx 3.5$, boosting the statistical density of collapsed massive objects by a factor of $4 \times 10^4$.

\section{Conclusion}
Unified Phase Dynamics resolves the JWST cosmic dawn crisis without Dark Matter particles, without super-Eddington accretion, and without fine-tuned top-heavy IMFs.

\begin{thebibliography}{99}
\bibitem{Carniani2024} S.~Carniani \textit{et al.}, Nature \textbf{633}, 318--322 (2024).
\bibitem{Labbe2023} I.~Labb{\'e} \textit{et al.}, Nature \textbf{616}, 266--269 (2023).
\bibitem{Finkelstein2023} S.~L.~Finkelstein \textit{et al.}, Astrophys. J. Lett. \textbf{946}, L13 (2023).
\bibitem{McGaugh2024} S.~S.~McGaugh, Astrophys. J. \textbf{966}, 210 (2024).
\bibitem{PavliukMonograph} H.~Pavliuk, CERN / Zenodo, Concept DOI: 10.5281/zenodo.24226129 (2026).
\bibitem{PavliukSPARC} H.~Pavliuk, CERN / Zenodo, Concept DOI: 10.5281/zenodo.24293023 (2026).
\bibitem{PavliukLensing} H.~Pavliuk, CERN / Zenodo, Concept DOI: 10.5281/zenodo.24293793 (2026).
\bibitem{PavliukCMB} H.~Pavliuk, CERN / Zenodo, Concept DOI: 10.5281/zenodo.24301701 (2026).
\end{thebibliography}

\end{document}
